% GENERATED FILE. Edit canonical lessons, then run npm run generate:books.
% canonical-source-hash: fnv1a64:7602155e81a39814
% canonical-lessons: ML-C340-iyite, ML-W340-ii, ML-C340-ennittu, ML-C340-atinusesham, ML-C340-gramathil, ML-R340-a-story-in-order

\chapter{Telling a Story in Order}
\label{ch:ml-telling-a-story-in-order}
\hlchaptermodality{340}

\begin{chapteropening}
\textbf{By the end of this chapter:} \emph{I can tell a short story about something that happened, in order, with ādyaṁ, enniṭṭŭ, atinuśēṣaṁ and oṭuvil, and read one of about a hundred words.}

Complete ML-R340-a-story-in-order: retell a day in order with enniṭṭŭ, -appōḷ and atinuśēṣaṁ, read the new letter \ml{ഈ}, and say what recently happened.
\end{chapteropening}

\section[\texorpdfstring{\ml{ഈയിടെ} (īyiṭe)}{ഈയിടെ (īyiṭe)}]{\ml{ഈയിടെ} (īyiṭe) — recently, lately}

\label{lesson:ML-C340-iyite}



\begin{quote}
Before the new one: say the Malayalam for \textquotedblleft{}meanwhile\textquotedblright{}, then for \textquotedblleft{}recently\textquotedblright{}. (\emph{itiniṭe}, \emph{aṭuttiṭe}.) Then \textquotedblleft{}once\textquotedblright{}. (\emph{orikkal}.)
\end{quote}

\subsection*{You'll want to know: \ml{ഈയിടെ}}
\textbf{\ml{ഈയിടെ}} — \emph{īyiṭe} — \textquotedblleft{}recently, lately\textquotedblright{}.

\begin{quote}
\textbf{\ml{ഈയിടെ} \ml{ഞാൻ} \ml{ഒരു} \ml{നല്ല} \ml{സിനിമ} \ml{കണ്ടു}.}
\end{quote}

\begin{quote}
\emph{īyiṭe ñān oru nalla sinima kaṇṭu.} — \textquotedblleft{}Recently I saw a good film.\textquotedblright{}
\end{quote}

With a past verb it means \textquotedblleft{}recently, a short while back\textquotedblright{}. With a present one it means \textquotedblleft{}lately, these days\textquotedblright{}:

\begin{quote}
\textbf{\ml{ഈയിടെ} \ml{മുത്തശ്ശിക്ക്} \ml{സുഖമില്ല}.}
\end{quote}

\begin{quote}
\emph{īyiṭe muttaśśikkŭ sukhamilla.} — \textquotedblleft{}Lately Grandmother has not been well.\textquotedblright{}
\end{quote}

\textbf{\ml{അടുത്തിടെ}}, which you know, says the same. You will hear both, and a story often opens on one of them.

Its first letter is one you have not written yet. The next lesson takes it out of this word and puts it in your hand.

\begin{cousinweb}
\textbf{\ml{ഈയിടെ}} (\emph{īyiṭe}) is native Dravidian, and it is read as \textquotedblleft{}in this interval\textquotedblright{}. \textbf{\ml{ഈ}}, \emph{ī}, is the \textquotedblleft{}this\textquotedblright{} that stands before a noun, from the near demonstrative base (DEDR 410). \textbf{\ml{ഇട}}, \emph{iṭa}, is \textquotedblleft{}interval, space between\textquotedblright{}, with Tamil \emph{iṭai} (DEDR 448). The \textbf{-\ml{എ}} on the end places it in time. The same \textbf{\ml{ഇട}} sits inside \textbf{\ml{ഇതിനിടെ}}, \textquotedblleft{}meanwhile\textquotedblright{}.
\end{cousinweb}

\subsection*{Your turn}
Say these aloud:

\begin{itemize}
\raggedright
  \item \emph{īyiṭe}, slowly, then at speed
  \item \textquotedblleft{}Recently I saw a good film\textquotedblright{} — \emph{īyiṭe ñān oru nalla sinima kaṇṭu}
  \item \textquotedblleft{}Lately Grandmother has not been well\textquotedblright{} — \emph{īyiṭe muttaśśikkŭ sukhamilla}
  \item the other word for \textquotedblleft{}recently\textquotedblright{} you already had — \emph{aṭuttiṭe}
\end{itemize}

\begin{tcolorbox}[breakable,colback=teal!4,colframe=teal!35!black,title={Before you move on}]
What does \textbf{\ml{ഈയിടെ}} mean? (\textbf{recently, lately.}) What is it, taken apart? (\textquotedblleft{}This\textquotedblright{} + \textquotedblleft{}interval\textquotedblright{}: \emph{\textbf{ī}} + \emph{\textbf{iṭa}}.) Which word for \textquotedblleft{}meanwhile\textquotedblright{} holds the same \textbf{\ml{ഇട}}? (\textbf{\ml{ഇതിനിടെ}}, \emph{itiniṭe}.)
\end{tcolorbox}
\section[\texorpdfstring{\ml{ഈ} (ī)}{ഈ (ī)}]{\ml{ഈ} (ī) — the letter that opens \ml{ഈയിടെ}}

\label{lesson:ML-W340-ii}



\begin{quote}
Before the new one: write \textbf{\ml{ഇ}} and \textbf{\ml{ീ}} once each, and name them. (Short \emph{i} at the start of a word; long \emph{ī} after a consonant, as in \textbf{\ml{വീട്}}.)

Now say \textbf{\ml{ഈയിടെ}} (\emph{īyiṭe}), \textquotedblleft{}recently\textquotedblright{}. You met it a moment ago. This lesson takes its first letter out and puts it in your hand.
\end{quote}

\subsection*{Script you'll notice: \ml{ഈ}}
\textbf{\ml{ഈ}} — \emph{ī}, a long \emph{ī}, at the start of a word.

After a consonant, long \emph{ī} is the sign \textbf{\ml{ീ}}: \textbf{\ml{വീട്}}, \emph{vīṭŭ}. At the start of a word there is no consonant to carry a sign, so the vowel stands as a letter of its own, \textbf{\ml{ഈ}}.

It is \textbf{\ml{ഇ}} with a tail added on the right: a small lobe, a notch, and a broader lobe that sweeps back along its foot. You have read the same tail on \textbf{\ml{ഊ}}, which is \textbf{\ml{ഉ}} with that tail.

On its own, \textbf{\ml{ഈ}} is also a word: \textquotedblleft{}this\textquotedblright{}, in front of a noun. \textbf{\ml{ഈ} \ml{വീട്}} is \textquotedblleft{}this house\textquotedblright{}.

\subsection*{Writing: \ml{ഈ}}
\textbf{\ml{ഈ}} is two parts side by side: the whole of \textbf{\ml{ഇ}} on the left, and the tail on the right. Write the \textbf{\ml{ഇ}} you already know first, then add the tail beside it: a small lobe, a notch, and a broader lobe that sweeps back along its foot.

Write \textbf{\ml{ഈ}} yourself, slowly, and larger than it is printed. Copy the printed shape.

\subsection*{Your turn}
\begin{itemize}
\raggedright
  \item \emph{Find:} \ml{ഈ} at the start of \textbf{\ml{ഈയിടെ}}
  \item \emph{Trace it:} \ml{ഈ} three times, saying \emph{ī} as you finish each one
  \item \emph{Write it:} \ml{ഇ}, then \ml{ഈ}, from memory
  \item \emph{Write it:} \textbf{\ml{ഈയിടെ}}
  \item \emph{Read it:} \textbf{\ml{ഈയിടെ}} aloud once more
\end{itemize}

\begin{tcolorbox}[breakable,colback=teal!4,colframe=teal!35!black,title={Before you move on}]
Which letter is this — \ml{ഈ}? (\textbf{ī}, long \emph{ī} at the start of a word.) Which word did it come from? (\textbf{\ml{ഈയിടെ}}, \emph{īyiṭe}, \textquotedblleft{}recently\textquotedblright{}.) What is the difference between \ml{ഇ} and \ml{ഈ}? (\textbf{The tail on the right.})
\end{tcolorbox}
\section[\texorpdfstring{\ml{എന്നിട്ട്} (enniṭṭŭ)}{എന്നിട്ട് (enniṭṭŭ)}]{\ml{എന്നിട്ട്} (enniṭṭŭ) — and then}

\label{lesson:ML-C340-ennittu}



\begin{quote}
Say \textquotedblleft{}having come\textquotedblright{}. (\textbf{\ml{വന്നിട്ട്}} (\emph{vanniṭṭŭ}).) Then the word for \textquotedblleft{}then, after that\textquotedblright{}. (\textbf{\ml{പിന്നെ}} (\emph{pinne}).)
\end{quote}

\begin{grammarlens}[title={Grammar Lens: one sentence, with -\ml{ഇട്ട്}}]
You already know the past form with \textbf{-\ml{ഇട്ട്}} on it: \textquotedblleft{}having done\textquotedblright{}. Put it inside a sentence and it says \textquotedblleft{}did this, and then\textquotedblright{}:

\begin{quote}
\textbf{\ml{ഞാൻ} \ml{ചായ} \ml{കുടിച്ചിട്ട്} \ml{ഓഫീസിൽ} \ml{പോയി}.}
\end{quote}

\begin{quote}
\emph{ñān cāya kuṭicciṭṭŭ ōphīsil pōyi.} — \textquotedblleft{}I drank tea and then went to the office.\textquotedblright{}
\end{quote}

Piece by piece: \textquotedblleft{}I, tea, having drunk, to the office went\textquotedblright{}. The first action is finished before the second begins.
\end{grammarlens}

\begin{grammarlens}[title={Grammar Lens: two sentences, with \ml{എന്നിട്ട്}}]
To keep the two actions as two sentences, start the second with \textbf{\ml{എന്നിട്ട്}}:

\begin{quote}
\textbf{\ml{ഞാൻ} \ml{ചായ} \ml{കുടിച്ചു}. \ml{എന്നിട്ട്} \ml{ഓഫീസിൽ} \ml{പോയി}.}
\end{quote}

\begin{quote}
\emph{ñān cāya kuṭiccu. enniṭṭŭ ōphīsil pōyi.} — \textquotedblleft{}I drank tea. And then I went to the office.\textquotedblright{}
\end{quote}

\textbf{\ml{എന്നിട്ട്}} is usually explained as the \textbf{-\ml{ഇട്ട്}} ending on the same root as \textbf{\ml{എന്ന്}}: \textquotedblleft{}having said so\textquotedblright{}, \textquotedblleft{}that done\textquotedblright{}. It ties the next action to the one before.

\textbf{\ml{പിന്നെ}} is the looser, everyday \textquotedblleft{}then, after that\textquotedblright{}, and it can also mean \textquotedblleft{}later\textquotedblright{}. \textbf{\ml{എന്നിട്ട്}} ties the next action to the one done a moment before.

A listener who wants the story to go on asks \textbf{\ml{എന്നിട്ടോ}?} (\emph{enniṭṭō?}), \textquotedblleft{}and then?\textquotedblright{}
\end{grammarlens}

\subsection*{Your turn}
Say these aloud:

\begin{itemize}
\raggedright
  \item \textquotedblleft{}I drank tea and then went to the office\textquotedblright{}, in one sentence, with \emph{kuṭicciṭṭŭ}
  \item the same as two sentences, with \emph{enniṭṭŭ}
  \item \textquotedblleft{}I saw a film. And then I went home.\textquotedblright{} — \emph{ñān sinima kaṇṭu. enniṭṭŭ vīṭṭil pōyi.}
  \item what a listener says to hear more — \emph{enniṭṭō?}
\end{itemize}

\begin{tcolorbox}[breakable,colback=teal!4,colframe=teal!35!black,title={Before you move on}]
\textquotedblleft{}And then\textquotedblright{}, at the start of a sentence? (\textbf{\ml{എന്നിട്ട്}}, \emph{enniṭṭŭ}.) The same idea inside one sentence? (The \textbf{-\ml{ഇട്ട്}} form: \emph{\textbf{kuṭicciṭṭŭ}}.) How is \textbf{\ml{പിന്നെ}} different? (\textbf{Looser: \textquotedblleft{}after that\textquotedblright{}, and it can also mean \textquotedblleft{}later\textquotedblright{}.})
\end{tcolorbox}
\section[\texorpdfstring{\ml{അതിനുശേഷം} (atinuśēṣaṁ)}{അതിനുശേഷം (atinuśēṣaṁ)}]{\ml{അതിനുശേഷം} (atinuśēṣaṁ) — after that: a story in order}

\label{lesson:ML-C340-atinusesham}



\begin{quote}
Say \textquotedblleft{}at first\textquotedblright{}, \textquotedblleft{}after\textquotedblright{}, and \textquotedblleft{}finally\textquotedblright{}. (\emph{ādyaṁ}, \emph{śēṣaṁ}, \emph{oṭuvil}.)
\end{quote}

\begin{grammarlens}[title={Grammar Lens: \ml{അതിനുശേഷം}}]
\textbf{\ml{ശേഷം}} stands after the thing it comes after, and that thing takes \textbf{-\ml{ന്}}. Make the thing \textbf{\ml{അത്}}, \textquotedblleft{}that\textquotedblright{}, and the result is a sentence link:

\begin{quote}
\ml{അതിന്} + \ml{ശേഷം} $\to$ \textbf{\ml{അതിനുശേഷം}} — \emph{atinuśēṣaṁ} — after that
\end{quote}

A story then hangs on four posts: a time to open it, \textbf{\ml{ആദ്യം}}, \textbf{\ml{പിന്നെ}} or \textbf{\ml{എന്നിട്ട്}} or \textbf{\ml{അതിനുശേഷം}} in the middle, and \textbf{\ml{ഒടുവിൽ}} to close.

\begin{quote}
\textbf{\ml{കഴിഞ്ഞ} \ml{ശനിയാഴ്ച} \ml{ഞങ്ങൾ} \ml{നഗരത്തിൽ} \ml{പോയി}. \ml{ആദ്യം} \ml{ചന്തയിൽ} \ml{പോയി}. \ml{അതിനുശേഷം} \ml{ഒരു} \ml{ചെറിയ} \ml{കടയിൽ} \ml{ചായ} \ml{കുടിച്ചു}. \ml{ഒടുവിൽ} \ml{രാത്രി} \ml{വീട്ടിൽ} \ml{എത്തി}.}
\end{quote}

\begin{quote}
\emph{kaḻiñña śaniyāḻca ñaṅṅaḷ nagarattil pōyi. ādyaṁ cantayil pōyi. atinuśēṣaṁ oru ceṟiya kaṭayil cāya kuṭiccu. oṭuvil rātri vīṭṭil etti.} — \textquotedblleft{}Last Saturday we went to the city. First we went to the market. After that we drank tea in a small shop. Finally we got home at night.\textquotedblright{}
\end{quote}
\end{grammarlens}

\begin{grammarlens}[title={Grammar Lens: when it happened, -\ml{അപ്പോൾ}}]
You already know \textbf{\ml{വരുമ്പോൾ}}, \textquotedblleft{}when he comes\textquotedblright{}. For \textquotedblleft{}when he came\textquotedblright{}, take the past form you know as \textbf{\ml{വന്ന}} and add \textbf{-\ml{പ്പോൾ}}: \textbf{\ml{വന്നപ്പോൾ}} (\emph{vannappōḷ}). In the same way \textbf{\ml{എത്തി}} gives \textbf{\ml{എത്തിയപ്പോൾ}} (\emph{ettiyappōḷ}), \textquotedblleft{}when (I) arrived\textquotedblright{}, and \textbf{\ml{കണ്ടു}} gives \textbf{\ml{കണ്ടപ്പോൾ}} (\emph{kaṇṭappōḷ}), \textquotedblleft{}when (I) saw\textquotedblright{}.

\begin{quote}
\textbf{\ml{ഞാൻ} \ml{എത്തിയപ്പോൾ} \ml{മുത്തശ്ശി} \ml{കാത്തിരിക്കുകയായിരുന്നു}.}
\end{quote}

\begin{quote}
\emph{ñān ettiyappōḷ muttaśśi kāttirikkukayāyirunnu.} — \textquotedblleft{}When I arrived, Grandmother was waiting.\textquotedblright{}
\end{quote}

The \textbf{-\ml{പ്പോൾ}} clause sets the moment; the main verb says what was going on in it, or what happened next.
\end{grammarlens}

\subsection*{Your turn}
Say these aloud:

\begin{itemize}
\raggedright
  \item the Saturday story, one sentence at a time
  \item \textquotedblleft{}when he came\textquotedblright{}, \textquotedblleft{}when I arrived\textquotedblright{}, \textquotedblleft{}when I saw\textquotedblright{} — \emph{vannappōḷ}, \emph{ettiyappōḷ}, \emph{kaṇṭappōḷ}
  \item \textquotedblleft{}When I arrived, Grandmother was waiting.\textquotedblright{}
  \item three things you did yesterday, with \emph{ādyaṁ}, \emph{atinuśēṣaṁ} and \emph{oṭuvil}
\end{itemize}

\begin{tcolorbox}[breakable,colback=teal!4,colframe=teal!35!black,title={Before you move on}]
\textquotedblleft{}After that\textquotedblright{}? (\textbf{\ml{അതിനുശേഷം}}: \textbf{\ml{അതിന്}} + \textbf{\ml{ശേഷം}}.) \textquotedblleft{}When I arrived\textquotedblright{}? (\textbf{\ml{എത്തിയപ്പോൾ}}, \emph{ettiyappōḷ}.) Which word closes a story? (\textbf{\ml{ഒടുവിൽ}}, \emph{oṭuvil}.)
\end{tcolorbox}
\section[(reading)]{(reading) — a visit to the village}

\label{lesson:ML-C340-gramathil}



\begin{quote}
This passage is about a hundred words long. Every word in it was taught before; one new word and a few new forms are glossed after it.
\end{quote}

\begin{tcolorbox}[breakable,skin=enhanced,colback=orange!6,colframe=orange!45!black,boxrule=0.5pt,arc=1mm,left=6pt,right=6pt,top=4pt,bottom=4pt,fonttitle=\bfseries,title={Read this}]
\begin{quote}
\ml{ഈയിടെ} \ml{ഞാൻ} \ml{മുത്തശ്ശിയുടെ} \ml{ഗ്രാമത്തിൽ} \ml{പോയി}. \ml{ആദ്യം} \ml{ഞാൻ} \ml{ഒരു} \ml{ബസിൽ} \ml{കയറി}. \ml{ബസിൽ} \ml{വലിയ} \ml{തിരക്കായിരുന്നു}. \ml{ഞാൻ} \ml{ഒരു} \ml{പുസ്തകം} \ml{വായിക്കുകയായിരുന്നു}. \ml{പെട്ടെന്ന്} \ml{ബസ്} \ml{നിർത്തി}. \ml{ഞാൻ} \ml{ഇറങ്ങി}, \ml{എന്നിട്ട്} \ml{പതുക്കെ} \ml{നടന്നു}.
\end{quote}

\begin{quote}
\ml{ഞാൻ} \ml{എത്തിയപ്പോൾ} \ml{മുത്തശ്ശി} \ml{വീട്ടിൽ} \ml{കാത്തിരിക്കുകയായിരുന്നു}. \ml{എന്നെ} \ml{കണ്ടപ്പോൾ} \ml{മുത്തശ്ശി} \ml{ചിരിച്ചു}. \ml{അതിനുശേഷം} \ml{ഞങ്ങൾ} \ml{ചോറും} \ml{മീനും} \ml{തിന്നു}. \ml{ഞാൻ} \ml{കുട്ടിയായിരുന്നപ്പോൾ} \ml{ഞങ്ങൾ} \ml{എല്ലാ} \ml{ദിവസവും} \ml{കുളത്തിൽ} \ml{കുളിക്കുമായിരുന്നു}. \ml{അന്ന്} \ml{ഗ്രാമത്തിൽ} \ml{ഒരു} \ml{ചെറിയ} \ml{കട} \ml{ഉണ്ടായിരുന്നു}. \ml{ഇപ്പോൾ} \ml{അവിടെ} \ml{ഒരു} \ml{വലിയ} \ml{കട} \ml{ഉണ്ട്}.
\end{quote}

\begin{quote}
\ml{ഒടുവിൽ} \ml{രാത്രി} \ml{ഞാൻ} \ml{നഗരത്തിലേക്ക്} \ml{മടങ്ങി}. \ml{അത്} \ml{ഒരു} \ml{നല്ല} \ml{ദിവസമായിരുന്നു}.
\end{quote}

Read it once without stopping. Do not translate as you go; let the sentences arrive.

Then read it again and sort every verb: scene, event, or habit.
\end{tcolorbox}

\begin{culture}
One word is new: \textbf{\ml{തിരക്ക്}}, \textquotedblleft{}a crowd, a rush\textquotedblright{}, from \textbf{\ml{തിരക്കുള്ള}}. The ending \textbf{-\ml{ലേക്ക്}} is \textquotedblleft{}to\textquotedblright{}: \emph{nagarattilēkkŭ}, \textquotedblleft{}to the city\textquotedblright{}.

The events, in order, are plain pasts of verbs you know (\emph{kayaṟi} from \emph{kayaṟuka}): \emph{pōyi}, \emph{kayaṟi}, \emph{nirtti}, \emph{iṟaṅṅi}, \emph{naṭannu}, \emph{ciriccu}, \emph{tinnu}, \emph{maṭaṅṅi}. The scene: \emph{vāyikkukayāyirunnu}, \emph{kāttirikkukayāyirunnu}. The habit: \emph{kuḷikkumāyirunnu}.
\end{culture}

\subsection*{Your turn}
Say these aloud:

\begin{itemize}
\raggedright
  \item what the writer was doing when the bus stopped — \emph{ñān oru pustakaṁ vāyikkukayāyirunnu}
  \item what Grandmother did when she saw the writer — \emph{muttaśśi ciriccu}
  \item what they used to do every day — \emph{kuḷattil kuḷikkumāyirunnu}
  \item the passage aloud, once, without stopping
\end{itemize}

\begin{tcolorbox}[breakable,colback=teal!4,colframe=teal!35!black,title={Before you move on}]
Why is \emph{kāttirikkukayāyirunnu} in the \emph{-ukayāyirunnu} past? (Grandmother \textbf{was waiting}: it is scene.)

Which word set the moment the writer arrived? (\emph{\textbf{ettiyappōḷ}}, \textquotedblleft{}when I arrived\textquotedblright{}.)
\end{tcolorbox}
\section[(review)]{Review — a story in order}

\label{lesson:ML-R340-a-story-in-order}



\begin{quote}
Say \textquotedblleft{}recently\textquotedblright{}. (\emph{īyiṭe}.) Then write its first letter, \textbf{\ml{ഈ}}, once from memory.
\end{quote}

\subsection*{Your turn}
\begin{itemize}
\raggedright
  \item \emph{Say it:} \textquotedblleft{}I drank tea. And then I went to the office.\textquotedblright{}
  \item \emph{Say it:} \textquotedblleft{}When I arrived, Grandmother was waiting.\textquotedblright{}
  \item \emph{Say it:} the Saturday story: the city, first the market, after that tea, finally home
  \item \emph{Say it:} \textquotedblleft{}Recently I saw a good film. And then I went home.\textquotedblright{}
  \item \emph{Say it:} from the reading, what used to happen in the village every day
  \item \emph{Write it:} \textbf{\ml{ഈയിടെ}}, then \textbf{\ml{ഈ} \ml{വീട്}}, \textquotedblleft{}this house\textquotedblright{}
\end{itemize}

\begin{tcolorbox}[breakable,colback=teal!4,colframe=teal!35!black,title={Before you move on}]
\textquotedblleft{}And then\textquotedblright{}, opening a sentence? (\textbf{\ml{എന്നിട്ട്}}.) \textquotedblleft{}After that\textquotedblright{}? (\textbf{\ml{അതിനുശേഷം}}.) \textquotedblleft{}When I arrived\textquotedblright{}? (\textbf{\ml{എത്തിയപ്പോൾ}}.) \textquotedblleft{}Recently\textquotedblright{}? (\textbf{\ml{ഈയിടെ}}, which opens with \textbf{\ml{ഈ}}.) In the reading, who was waiting, and what was the writer doing on the bus? (\textbf{Grandmother; reading a book.})
\end{tcolorbox}
